\section*{Chapitre \ref{ch:pl:the-risk-grid} --- La grille de risque}

\begin{solution}{exo:pl:the-risk-grid:1}
2\,200 tâches (220 lots de transactions multipliés par dix lots de scénarios), soit $2\,200 \times 5 = 11\,000$ secondes-cœur de coût fixe, environ 10~\% de la
valorisation.
\end{solution}

\begin{solution}{exo:pl:the-risk-grid:2}
La fin est fixée par une seule tâche --- le lot de la transaction lente sur tous les scénarios, environ 5\,380~secondes --- et par le moment où elle démarre. Davantage de cœurs
la font démarrer plus tôt, puisque la file d'attente qui la précède se vide plus vite, mais ne peuvent pas la raccourcir~: sur 64 cœurs elle se termine encore 6\,922~secondes après le
début.
\end{solution}

\begin{solution}{exo:pl:the-risk-grid:3}
Seize revalorisations (chacun des huit piliers à la hausse et à la baisse) plus la base~: 17 appels~; l'adjoint enregistre une valorisation et exécute un balayage inverse,
pour un coût de quelques valorisations, quel que soit le nombre de piliers.
\end{solution}

\begin{solution}{exo:pl:the-risk-grid:4}
Le nombre de tâches croît à mesure que les lots rétrécissent --- 44\,000 pour cinq scénarios --- et chacune porte cinq secondes de coût fixe~; en dessous
d'environ 50 scénarios, les coûts fixes croissent plus vite que l'équilibre ne s'améliore.
\end{solution}

\begin{solution}{exo:pl:the-risk-grid:5}
La swaption défaillante se trouve dans un lot de plusieurs milliers de transactions~; chaque tâche de ce lot --- une par lot de scénarios --- échoue, si bien que les
résultats du lot, pour chaque scénario, sont perdus~: 38.3~millions de cellules.
\end{solution}

\begin{solution}{exo:pl:the-risk-grid:6}
Lorsque la version des données de marché change, puisque la clé de chaque résultat change. Le risque intrajournalier sur un nouveau marché est couvert par les sensibilités
reportées (chapitre~18) et, si besoin, une réévaluation partielle des transactions les plus exposées au mouvement.
\end{solution}

\begin{solution}{exo:pl:the-risk-grid:7}
Sur 64 cœurs, le plan naïf respecte la fenêtre de 5 à 500 scénarios par tâche (meilleur cas à 50~: 2\,319~secondes) et n'échoue qu'à 1\,000~; le
plan tenant compte des coûts le respecte à toute granularité, de 5 à 1\,000 (meilleur cas à 250~: 1\,989~secondes). Davantage de cœurs élargissent la plage favorable~; ils ne suppriment pas la nécessité de planifier.
\end{solution}

\begin{solution}{exo:pl:the-risk-grid:8}
Doubler les cœurs a fait passer la fin de 06~h~51 à 06~h~25, toujours tardive, car le retard venait d'une seule tâche et du plan, et non de la capacité.
Mesurer les tâches, isoler ce qui est lent, choisir la granularité~; n'acheter de la capacité que lorsque la borne inférieure d'un bon plan dépasse la fenêtre.
\end{solution}

\begin{solution}{pb:pl:the-risk-grid:1}
\begin{enumerate}
\item 105\,000 transactions par 1\,000 scénarios~: 105~millions d'appels de valorisation (auxquels s'ajoutent les valorisations de base).
\item Les 5\,000 autocallables Monte-Carlo~: 22.2~millisecondes par appel, 99~\% des 112.2~secondes de valorisation par scénario.
\item Cinq secondes~: chargement de l'instantané et des transactions, écriture des résultats.
\item 90~minutes, de 04~h~30 à 06~h~00, sur 32 cœurs.
\item Environ 4\,000~secondes à 100 scénarios par tâche~: le travail réparti sur 32 cœurs.
\item 2.35, 1.19 et 2.93~heures.
\item 20, 50, 100 et 250 scénarios par tâche (20 par seize secondes).
\item Il crée une tâche de plus de 5\,300~secondes, lancée tard parce que son estimation est faible.
\item Il utilise le coût mesuré la nuit précédente pour donner aux transactions lentes leurs propres tâches, réparties finement sur les scénarios.
\item 05~h~34, avec 250 scénarios par tâche.
\item Par tâche~: deux nouvelles tentatives, puis les résultats de la tâche sont perdus~; par transaction~: la transaction est mise en quarantaine avec un rapport et la tâche se termine.
\item 38.3~millions de résultats et 1\,100 secondes-cœur par tâche~; 1\,000 résultats et rien par transaction.
\item 5.9 fois (388 contre 66~microsecondes).
\item 1.02~millions de revalorisations par chocs remplacées par 60\,000 valorisations enregistrées.
\item 1.5~millions d'appels pour 1\,500 transactions modifiées, contre 105~millions.
\item \textbf{Résultat principal.} Sur 32 cœurs, la grille prend 2.35~heures avec une tâche par lot de transactions sur tous les scénarios, 1.19 avec 100 scénarios
par tâche et 2.93 avec cinq~; seules les tailles de 20 à 250 scénarios respectent la fenêtre de 90 minutes avec le plan naïf, et un plan tenant compte des coûts la respecte à toute
granularité de 20 à 1\,000 (meilleur cas à 250~: 1.07~heures)~; doubler les cœurs fait passer le plan naïf de 06~h~51 à 06~h~25~; les adjoints remplacent 1.02~millions
de revalorisations par chocs des swaps par 60\,000 valorisations, et une relance incrémentale intrajournalière fait 1.5~millions d'appels au lieu de 105~millions.
\item Planifier selon le coût mesuré, le lot de scénarios étant choisi en fonction du coût fixe et du cluster, et mettre en quarantaine les transactions défaillantes.
\item La durée de chaque tâche comparée à son estimation, les transactions les plus lentes, les échecs et les mises en quarantaine, la fin comparée à la fenêtre, et la borne
inférieure.
\item Lorsque la borne inférieure d'un bon plan, et non sa \omterm{def:pl:distributed-compute-and-schedulers:makespan}{durée totale} (makespan), approche de la fenêtre.
\item Ses tâches les plus longues et le moment où elles démarrent.
\end{enumerate}
\end{solution}

\begin{solution}{iq:pl:the-risk-grid:1}
Les tâches du chemin critique~: lesquelles ont fini en dernier et ont duré le plus longtemps, et si leurs durées correspondaient à leurs estimations~; viennent ensuite les échecs et les nouvelles tentatives~;
puis la granularité au regard du coût fixe d'une tâche~; la capacité en dernier.
\iqlookfor{Chemin critique avant capacité.}
\end{solution}

\begin{solution}{iq:pl:the-risk-grid:2}
Regrouper les transactions en lots selon le coût mesuré pour viser une cible par scénario, regrouper les scénarios de sorte que les tâches durent quelques minutes (plusieurs fois le coût fixe, bien moins que
la fenêtre), isoler et répartir les transactions coûteuses, et planifier les plus longues en premier.
\iqlookfor{Regroupement selon le coût, deux dimensions.}
\end{solution}

\begin{solution}{iq:pl:the-risk-grid:3}
Parce qu'il fournit chaque sensibilité d'un prix pour quelques fois le coût d'un seul prix, alors que la méthode par chocs coûte deux prix par sensibilité~; il n'aide pas
lorsqu'il y a peu de sensibilités, lorsque le code ne peut pas être enregistré sur bande (un pricer à code fermé), ou pour des payoffs non lisses sans lissage.
\iqlookfor{Coût indépendant du nombre d'entrées.}
\end{solution}

\begin{solution}{iq:pl:the-risk-grid:4}
La mettre en quarantaine~: la boucle de valorisation intercepte son échec, l'inscrit dans un rapport d'erreur avec sa cause, achève la tâche sans elle, et le
rapport est transmis au responsable des données de la transaction avant le prochain passage.
\iqlookfor{Isoler l'échec, jamais la tâche.}
\end{solution}

\begin{solution}{iq:pl:the-risk-grid:5}
Une grille planifiée à partir des coûts mesurés sur un ordonnanceur à partage équitable et à ordre des plus longues en premier~; des tâches dimensionnées d'après leurs coûts fixes~; des résultats
stockés par transaction, version de la transaction et version des données de marché~; échecs mis en quarantaine avec leurs rapports~; sensibilités adjointes lorsque la bibliothèque les prend en charge~;
\omterm{def:pl:the-risk-grid:incremental}{réévaluations incrémentales} intrajournalières pour les transactions modifiées, et sensibilités reportées pour les mouvements de marché.
\iqlookfor{Planification, robustesse, réutilisation.}
\end{solution}
